\subsection{The four-sphere minimum and explicit three-dimensional models}
\begin{theorem}\label{thm:C-S4-nine}
The minimum number of arrows of a finite category of type $S^4$ is $9$.
\end{theorem}
\begin{proof}
The nine-element monoid and complete finite projective resolution in
Proposition~\ref{prop:B-S4-nine} prove the upper bound.
For monoids of smaller order, repeatedly apply the one-sided corner
reductions used in the census. Among all the irreducible monoids of
orders two through seven there is no table with
$H_1=H_2=H_3=0$; the lists contain respectively $1,1,2,2,7,25$
isomorphism-or-opposite classes, and their exact integral homology
is recorded in \texttt{c\_monoid\_homology.json}.
For a remaining nongroup irreducible monoid of order eight, the
minimal ideal is a Rees matrix semigroup $I\times H\times J$.
If $|I|=1$ or $|J|=1$, its maximal subgroup is a reducing corner,
contradicting irreducibility. Thus $|I|,|J|\ge2$.
Since the identity lies outside this proper ideal, its size is at
most seven. Consequently $H=1$ and, up to opposition, its dimensions
are $2\times2$ or $2\times3$.
The independently generated constrained tables give $243$ and $8$
irreducible isomorphism-or-opposite classes, respectively.
Here the raw completion counts with the band fixed in its stated
labels are $22746$ and $56$; these are not counts of all labelled
order-eight monoids.

The programs \texttt{c\_order8\_filter.py} and
\texttt{c\_order8\_survivors.py} exclude $244$ tables by mod-two
homology in degrees one through three. Each of the seven remaining
tables has $H_4(-;\F_2)=0$, also excluding $S^4$.
Groups are excluded by their nontrivial universal groups and
Section~\ref{sec:A}. This is a complete structural enumeration of
the order-eight irreducible cases, not an inference from a randomly
chosen family.

For categories with several objects the complete eight-arrow
enumeration has the counts in Proposition~\ref{prop:C-census}.
After exact isomorphism and opposition checks, the only candidate
with mod-two Betti numbers $(1,0,0,0,1)$ through degree four is
the categorical join of the five-element sphere monoid $P_2$ and
the two-element group $C_2$. A categorical join adds one arrow
from each object of the first category to each object of the second,
with its forced compositions. Thus this category has
$5+2+1=8$ arrows. Its nerve realizes the topological join, so
\[
 B(P_2*C_2)\simeq S^2*\RP^\infty
 \simeq\Sigma^3\RP^\infty.
\]
Its fourth integral homology is $\Z/2$, not $\Z$; the independent
integer Smith calculation in \texttt{c\_categories8\_residual.py}
also checks this directly.

This exact eight-arrow test covers all smaller multiobject categories
as well. To see the reduction, one may add a new identity at any
one object $x$, leaving its old identity as an idempotent $e$ that
acts as the identity on all old arrows. This adds exactly one arrow.
In the Karoubi envelope the full subcategory on $(x,e)$ and the other
original objects is the old category. The functor sending $x$ to
$(x,e)$ and every old arrow to itself is naturally related to the
canonical inclusion by the component $e$ at $x$ and identities
elsewhere. These functors prove that the added arrow preserves
the classifying-space homotopy type. Iterate to reach eight arrows.
A connected category with at most eight arrows has at most four
objects, because it needs at least $k-1$ nonidentity arrows for
$k$ objects. Hence every case reached is in the displayed census.
\end{proof}

\begin{proposition}\label{prop:C-lens-uppers}
There are explicit finite categories with four objects and $16$
arrows of type $\RP^3$, and with four objects and $20$ arrows of
type $L(3,1)$.
\end{proposition}
\begin{proof}
For the first use objects $0<1<2<3$, take
$\operatorname{Hom}(i,j)=C_2$ for every $i<j$, and compose by
addition; all endomorphisms are identities. Its covering poset
has objects $(i,g)$, $g\in C_2$, and a unique arrow for every
strict increase in $i$. Its order complex is the join of four
two-point sets, the cross-polytope boundary $S^3$.
Simultaneously adding $1$ to the binary coordinates is the antipodal
action. The orbit category is the stated category and each nerve
chain lifts uniquely after its first vertex is chosen. Therefore
its classifying space is homeomorphic to $S^3/C_2=\RP^3$.
The arrow count is $4+2\binom42=16$.

For the second again use objects $0<1<2<3$. In $C_3$ set
\[
 \operatorname{Hom}(0,1)=\operatorname{Hom}(2,3)=\{0,1\},
 \qquad \operatorname{Hom}(i,j)=C_3
 \quad\text{for the four other pairs }i<j,
\]
and compose by addition. The covering poset on $(i,g)$ has an
arrow precisely when the difference of its group coordinates lies
in the corresponding displayed hom-set. The subposet on levels
$0,1$ is the face poset of a triangular circle, as is the subposet
on levels $2,3$; every object of the first is less than every
object of the second. Its order complex is consequently
$S^1*S^1\cong S^3$. The deck generator rotates both circles
by the same $2\pi/3$ angle. In the usual join identification with
the unit sphere of $\mathbb C^2$, this is
$(z_1,z_2)\mapsto(\zeta_3z_1,\zeta_3z_2)$.
Its quotient is the lens space $L(3,1)$, with this explicit
convention for the lens action. The arrow count is
$4+2+2+4\cdot3=20$.

The program \texttt{c\_lens\_models.py} writes both complete typed
composition tables, checks every unit and associativity equation,
constructs and checks both covering posets and their unique lifts,
and computes integral homology through their full dimension three.
The geometric covering identifications prove the asserted homotopy
types; the homology check is an independent certificate.
The lower bounds for these two models are addressed separately;
the present statement asserts only the two explicit upper bounds.
\end{proof}
